%==============================================================================
\section{Rendezvous Server}
\label{sec:rendezvous}
%==============================================================================
A \emph{rendezvous server} is a separate process---a host that provides a static, globally-routable IP address and runs a lightweight GLP \emph{rendezvous agent}.  Its purpose is to give NAT-bound agents, which cannot otherwise reach one another, a fixed point they can always reach outbound, from which they discover their own public addresses and coordinate UDP hole-punching.

This section specifies the rendezvous server only by its networking requirements---what a machine must provide to host one.  The rendezvous agent's coordination behaviour (matching reconnection and availability requests, relaying addresses, verifying capabilities) is GLP agent code, built on the system predicates of Section~\ref{sec:system-predicates}, and is out of scope here; so is the reconnection protocol that uses it.

\paragraph{Networking requirements.}
\begin{itemize}
\item \textbf{Static, globally-routable address.}  The host has a fixed public \code{ip:port} that any agent can reach outbound, traversing NAT in the outbound direction without assistance.
\item \textbf{Address observation.}  When an agent connects, the rendezvous server observes the source \code{ip:port} of the agent's packets and exposes it via \code{peer\_address/2} (Section~\ref{sec:system-predicates}).  This is the address-discovery mechanism that lets a NAT-bound agent learn its own public address; no external STUN infrastructure is required.
\item \textbf{Hole-punch coordination.}  The rendezvous server relays observed addresses between two agents and, via \code{punch\_udp/1}, has each send punch traffic toward the other so a direct UDP path can form.  Once the path is open the peers run UDX and Noise~XX directly (Sections~\ref{sec:ip-connection} and~\ref{sec:session}); the rendezvous server sees signaling metadata and observed addresses, never the encrypted payload.
\item \textbf{Capability verification.}  Before coordinating, the rendezvous agent verifies the signed capability a connecting agent presents---a reconnection attestation (a signed binding of two public keys), or an invite capability for a first-contact cold call (Section~\ref{sec:ip-cold-call})---via \code{verify\_attestation/4}.  The networking layer checks the signature; what the binding authorizes is the rendezvous agent's GLP logic.
\item \textbf{No relay.}  The rendezvous server is not a TURN relay; it never forwards madGLP payloads.  When direct connectivity cannot be established (e.g., symmetric NAT), relaying is an application-level concern handled in GLP by a connected common peer, not by the networking layer.
\item \textbf{Not a social-graph member.}  A rendezvous agent has no friends list and does not participate in the social graph; it accepts connections from any agent.
\end{itemize}

\paragraph{Federation.}
Any individual or organization can run a rendezvous server.  The architecture is federated: agents may use several rendezvous servers for redundancy, and the system operates with no rendezvous server at all when agents are BLE-adjacent (Section~\ref{sec:ble}).  In the terminology of the ICE framework~\cite{rfc8445}, a rendezvous server combines the roles of STUN server (address discovery via \code{peer\_address/2}) and application-specific signaling server (candidate exchange and hole-punch coordination, performed in GLP); it does not serve as a TURN relay.
